\documentclass[11pt,a4paper]{article}
\usepackage[T1]{fontenc}
\usepackage{lmodern,microtype}
\usepackage[margin=26mm,headheight=15pt]{geometry}
\usepackage{amsmath,amssymb,amsthm,mathtools}
\usepackage{booktabs,longtable,array,enumitem}
\usepackage[dvipsnames]{xcolor}
\usepackage{fancyhdr}
\usepackage[colorlinks=true,linkcolor=MidnightBlue,citecolor=MidnightBlue,urlcolor=MidnightBlue,breaklinks=true]{hyperref}
\usepackage{bookmark}
\pagestyle{fancy}\fancyhf{}
\lhead{\small Reciprocal Lerch convolutions}\rhead{\small ProveIt research continuation}
\cfoot{\thepage}\renewcommand{\headrulewidth}{0.3pt}
\numberwithin{equation}{section}
\newtheorem{theorem}{Theorem}[section]
\newtheorem{lemma}[theorem]{Lemma}
\newtheorem{proposition}[theorem]{Proposition}
\newtheorem{corollary}[theorem]{Corollary}
\theoremstyle{definition}\newtheorem{definition}[theorem]{Definition}
\theoremstyle{remark}\newtheorem{remark}[theorem]{Remark}
\DeclareMathOperator{\Li}{Li}
\DeclareMathOperator{\Rea}{Re}
\DeclareMathOperator{\Ima}{Im}
\newcommand{\dd}{\,\mathrm d}
\newcommand{\C}{\mathbb C}\newcommand{\R}{\mathbb R}\newcommand{\N}{\mathbb N}
\newcommand{\E}{\mathcal E}
\newcommand{\CC}{\mathcal C}
\newcommand{\MM}{\mathcal M}
\newcommand{\stir}[2]{\genfrac{[}{]}{0pt}{}{#1}{#2}}
\newcommand{\src}[1]{\nolinkurl{#1}}
\newcommand{\pin}{\src{0c740647c48baec7f0b4192723e977d69917dd52}}
\setlist{itemsep=4pt,topsep=6pt}
\allowdisplaybreaks[2]
\hypersetup{pdftitle={Reciprocal Lerch Convolutions},pdfauthor={Research continuation for Vladimir Reshetnikov's ProveIt programme},pdfsubject={Entire-order identities, central-factorial reductions, logarithmic moments, and Stieltjes--Gamma calculus}}
\title{\textbf{Reciprocal Lerch Convolutions}\\[6pt]
\large Entire-order identities, central-factorial zeta reductions,\\
all logarithmic moments, and Stieltjes--Gamma calculus}
\author{Research continuation for Vladimir Reshetnikov's ProveIt programme\\
\small Mathematical derivations and computational checks with OpenAI assistance}
\date{10 October 2026}
\begin{document}
\maketitle
\begin{abstract}
A reciprocal-argument product of two Lerch profiles has an ordinary
integral equal to $(s+t)\zeta(s+t+1,a)$, interpreted as an entire function
of the total order. We develop this order-addition mechanism at every
multiplicative depth. An explicit central-factorial polynomial reduces
an $r$-fold convolution to finitely many Lerch values, and at unit
fugacity to Hurwitz-zeta or alternating Hurwitz-zeta values. A separate
finite differential-polynomial calculus evaluates \emph{every} mixed
logarithmic moment for arbitrary complex orders. For two factors we
obtain a closed formula whose apparent zeta poles cancel by a polynomial
divisibility identity.

The resulting order derivatives give ordinary convergent products
representing generalized Stieltjes constants. Higher-depth derivatives
involve finite Stirling and generalized-harmonic polynomials, positive
odd zeta jets, and an explicit Stieltjes term. Gamma-weighted contour
integrals, reciprocal polylogarithmic remainders, a Gaussian odd-part
family, unequal-shift partial fractions, and anchored antiderivatives
are derived in the same normalization. Negative-order resonances retain
nonzero residue contributions; one explicit depth-four term is $1/3$.

The article provides full analytic proofs, exact finite checks, numerical
diagnostics, and integration notes. It continues the repository's
Mellin and Stieltjes programme in a distinct reciprocal geometry; it
does not claim to settle the same-direction two-scale problem, the
Gaussian $S_6/S_8$ conjectures, or historical priority for the underlying
convolution method.
\end{abstract}
\vfill
\noindent\small\textbf{Source anchor.} Targeted repository reading is pinned to\\
\pin. Three accessible incoming packages were inspected in relevant
parts; three newer binary packages were listed but not inspected.
No repository files were modified.\normalsize
\clearpage
\begingroup
\small
\tableofcontents
\endgroup
\clearpage
\section{Research target, contribution, and provenance}
\label{rlc:sec:scope}

The canonical ProveIt polylogarithm manuscript contains an extensive theory of harmonic sums, cyclotomic
polylogarithms, Hurwitz order derivatives, Stieltjes distributions,
and Gamma correlations \cite{Repo}. Its Mellin--Lerch continuation asks
for exact product integrals and their spectral derivatives \cite{MLD}.
The incoming bilinear report treats positive-integer polylogarithm
orders under same-direction rational Mellin kernels, with a finite
multiple-polylogarithm closure theorem and further questions about
arbitrary orders \cite{BMH}.

Here the change in geometry is essential. We study reciprocal arguments
and, more generally, multiplicative constraints
$y_1\cdots y_r=e^\mu$. This is not a relabeling of the same-direction
integral
\begin{equation}\label{rlc:eq:priorproblem}
 \int_0^\infty \frac{x^{b-1}\Li_s(-zx)\Li_t(-wx)}{(1+x)^N}\dd x.
\end{equation}
The reciprocal geometry has a simple bilateral transform, and that
transform permits an entire-order finite closure absent from the
integer-only cone expansion of the preceding report. Our results
solve the explicitly formulated reciprocal problem, not all of
\eqref{rlc:eq:priorproblem}.

\subsection{A first identity and the central normalization}
Define, for $y>0$ and $\Rea a>0$,
\begin{equation}\label{rlc:eq:profile}
 F_{a,s}(y)=y\Phi(-y,s,a),\qquad
 f_{a,s}(x)=F_{a,s}(e^x),\qquad
 \E_a(v)=v\zeta(v+1,a).
\end{equation}
The last function is filled in at $v=0$ by $\E_a(0)=1$; it is entire
in $v$. Our normalization gives
\begin{equation}\label{rlc:eq:specialprofiles}
 F_{1,s}(y)=-\Li_s(-y),\qquad
 F_{a,0}(y)=\frac{y}{1+y}.
\end{equation}
The most economical instance of the theory is
\begin{equation}\label{rlc:eq:opening}
 \boxed{\displaystyle
 \int_0^\infty F_{a,s}(y)F_{a,t}(1/y)\frac{\dd y}{y}
      =\E_a(s+t),\qquad s,t\in\C.}
\end{equation}
All the displayed integrals in this article, unless explicitly described
as contour integrals, are ordinary absolutely convergent integrals.
No distribution product or finite-part cutoff is hidden in
\eqref{rlc:eq:opening}. In particular, the value at $s+t=0$ is $1$, not
zero and not a directional limit.

The full results go considerably beyond this two-factor identity.
Theorems~\ref{rlc:thm:addition} and \ref{rlc:thm:central} treat every
number of factors and every complex order. Theorem~\ref{rlc:thm:allmoments}
evaluates all mixed powers of their logarithmic arguments by a finite
recurrence. Theorem~\ref{rlc:thm:twomoments} gives a particularly short
closed two-factor formula. The Stieltjes, harmonic, Gamma, and
antiderivative consequences are proved in
Sections~\ref{rlc:sec:stieltjes}--\ref{rlc:sec:primitives}.

\subsection{Conventions and classical inputs}
We use
\[
 \Phi(z,s,a)=\sum_{n=0}^\infty\frac{z^n}{(n+a)^s},\qquad
 \zeta(s,a)=\sum_{n=0}^\infty(n+a)^{-s}
\]
in their initial domains and their standard analytic continuations.
For $\Rea a>0$, $\log(n+a)$ is the branch in the right half-plane.
Pochhammer symbols are rising factorials:
$(s)_0=1$, $(s)_j=s(s+1)\cdots(s+j-1)$.
A superscript in square brackets denotes an order derivative, while
$\psi^{(k)}$ denotes an argument derivative of the digamma function.
Generalized Stieltjes constants have the convention
\begin{equation}\label{rlc:eq:stieltjesconvention}
 \zeta(1+v,a)=\frac1v+
 \sum_{n=0}^\infty\frac{(-1)^n}{n!}\gamma_n(a)v^n;
 \qquad \gamma_0(a)=-\psi(a).
\end{equation}
These conventions agree with the standard Hurwitz, Lerch, Fermi-integral,
and Gamma identities summarized in the DLMF \cite{DLMF}.

The elementary mechanisms are classical: Euler's beta integral,
bilateral convolution, the semigroup law for one-sided fractional
integrals, and partial fractions. Real convolution methods for
Fermi--Dirac-type integrals already appear explicitly in
Selvaggi and Selvaggi \cite{Selvaggi}. The incoming unequal-twist report
also uses resolvent partial fractions, harmonic coefficients, and
monodromy, in a periodic Fourier setting \cite{UTJ}.
We do not present those methods as discoveries. The contribution here
is their explicit entire-order, all-depth, all-logarithmic-moment
specialization and the exact identities it produces. We make no claim
that the initial two-factor integral, or every special case below, is
new in the wider literature.

\subsection{What the source audit does and does not establish}
The source pin was read through the connected GitHub repository. The
canonical README, the Gaussian quotient chapter, and the Mellin
research-question section were inspected at that pin. Relevant sections
of the incoming bilinear, unequal-twist, and nonlinear-transport packages
were also available through the user's Library. Three further incoming packages, named for cubic harmonic mixed orders,
exact resonant identities, and independent orders, were listed in the
incoming directory but were not available for content inspection. Their
complete filenames are recorded in the source-audit artifact.
Their names alone establish neither overlap nor mathematical content.

The existing $S_4$ proof is not counted again. The current short Gaussian
$S_6$ reduction and revised $S_8$ candidate remain conjectural in the
inspected canonical manuscript. Our Gaussian specialization is an
ordinary integral family; it is not a zero certificate for either
frozen relation vector. No proof-assistant formalization is supplied.

\section{Analytic profiles and their bilateral transform}
\label{rlc:sec:analysis}

Write $f_0(x)=e^x/(1+e^x)$. For $\Rea s>0$ the Lerch integral gives
\begin{equation}\label{rlc:eq:fractional}
 f_{a,s}(x)=\frac1{\Gamma(s)}
 \int_0^\infty u^{s-1}e^{-(a-1)u}f_0(x-u)\dd u.
\end{equation}
The right side explains the covariant derivative
$D_a=\partial_x+a-1$:
\begin{equation}\label{rlc:eq:covariant}
 D_a f_{a,s}=f_{a,s-1}.
\end{equation}
The analytic justification of all subsequent parameter operations is
concentrated in the next lemma.

\begin{lemma}[Weighted holomorphy]\label{rlc:lem:weighted}
Fix $\Rea a>0$ and choose a real $c$ satisfying
\begin{equation}\label{rlc:eq:cstrip}
 \max(0,1-\Rea a)<c<1.
\end{equation}
Then $s\mapsto e^{-cx}f_{a,s}(x)$ is an entire family in
$L^1(\R)\cap L^\infty(\R)$. The same is true after multiplication by
any fixed power of $x$, after any fixed order derivative, and locally
jointly in $a$. The profile is smooth in $x$, and
\eqref{rlc:eq:covariant} holds for all $s\in\C$.
\end{lemma}
\begin{proof}
Every derivative of $f_0$ is $f_0$ times a polynomial in $f_0$, so
$|f_0^{(j)}(x)|\le C_j f_0(x)$ on the real line. The function
$e^{-cx}f_0(x)$ is bounded and integrable for $0<c<1$.
For $\Rea s>0$, multiply \eqref{rlc:eq:fractional} by $e^{-cx}$ and
put $v=x-u$. Its absolute integral is bounded by a constant times
\[
 \int_0^\infty u^{\Rea s-1}
 e^{-(c+\Rea a-1)u}\dd u
 \int_\R e^{-cv}f_0(v)\dd v.
\]
Both factors are finite. The same estimate with the second factor a
supremum proves boundedness. On compact subsets of the parameter
half-planes, derivatives insert only fixed powers of $u$ and $\log u$,
so the estimates prove holomorphy in these norms.

For $\Rea s>1$, differentiation and integration by parts in $u$ give
$D_a f_{a,s}=f_{a,s-1}$: the derivative of
$e^{-(a-1)u}f_0(x-u)$ is its negative covariant $x$ derivative, and
both boundary terms vanish. For a compact set of arbitrary complex
$s$, choose an integer $N$ with $\Rea(s+N)>0$ and define
$f_{a,s}=D_a^Nf_{a,s+N}$. The integration-by-parts identity makes these
definitions consistent on overlaps. They agree with the original Lerch
series when $x<0$ and hence with its continuation. The estimates for
$f_0^{(j)}$ prove the asserted norm holomorphy everywhere.

Finally choose $\delta>0$ so that $c-\delta$ and $c+\delta$ both satisfy
\eqref{rlc:eq:cstrip}. The elementary inequality
$|x|^m\le C_{m,\delta}e^{\delta|x|}$ transfers the two neighboring
weighted estimates to every power of $x$. Cauchy's formula in $s$ and
$a$ gives the same conclusions for their derivatives.
\end{proof}

\begin{theorem}[Bilateral transform]\label{rlc:thm:transform}
For $s\in\C$, $\Rea a>0$, and
$\max(0,1-\Rea a)<\Rea p<1$,
\begin{equation}\label{rlc:eq:transform}
 \int_\R e^{-px}f_{a,s}(x)\dd x
 =\Gamma(p)\Gamma(1-p)(p+a-1)^{-s}
 =\frac{\pi}{\sin\pi p}(p+a-1)^{-s}.
\end{equation}
The power uses the right-half-plane logarithm of $p+a-1$.
\end{theorem}
\begin{proof}
For $\Rea s>0$, Fubini applies to \eqref{rlc:eq:fractional} by the
preceding estimates. The substitution $x=v+u$ separates the transform
into
\[
 \frac1{\Gamma(s)}\int_0^\infty
 u^{s-1}e^{-(p+a-1)u}\dd u
 \int_\R e^{-pv}f_0(v)\dd v.
\]
The first factor is $(p+a-1)^{-s}$. With $y=e^v$, the second is
$\int_0^\infty y^{-p}/(1+y)\dd y=B(1-p,p)$.
Euler's reflection formula gives the stated Gamma product. Both sides
are entire in $s$, so the identity theorem completes the proof.
\end{proof}

There is no factor $\Gamma(s)$ left in \eqref{rlc:eq:transform}.
Losing this normalization would destroy order addition and every
Stieltjes specialization. The common strip depends on the shift $a$,
not on the spectral order $s$. This distinction is what makes the
all-complex-order statements possible.

\section{Order addition at every multiplicative depth}
\label{rlc:sec:convolution}

For $r\ge2$, set $x_r=\mu-x_1-\cdots-x_{r-1}$ and define
\begin{equation}\label{rlc:eq:convolution}
 (f_{a,s_1}*\cdots*f_{a,s_r})(\mu)
 =\int_{\R^{r-1}}\prod_{j=1}^r f_{a,s_j}(x_j)
       \dd x_1\cdots\dd x_{r-1}.
\end{equation}
Equivalently $y_j=e^{x_j}$ and $y_1\cdots y_r=e^\mu$, with measure
$\prod_{j<r}\dd y_j/y_j$. For $r=1$ use ordinary evaluation.

\begin{theorem}[Entire-order addition]\label{rlc:thm:addition}
For $\Rea a>0$, real $\mu$, and arbitrary $s_1,\ldots,s_r\in\C$,
\eqref{rlc:eq:convolution} is absolutely convergent and depends on the
orders only through $W=s_1+\cdots+s_r$. Denote it by $\CC_{r,a}(W;\mu)$.
It is entire in $W$ and has the representation
\begin{equation}\label{rlc:eq:barnes}
 \CC_{r,a}(W;\mu)=\frac1{2\pi i}\int_{c-i\infty}^{c+i\infty}
 e^{p\mu}\left(\frac{\pi}{\sin\pi p}\right)^r
 (p+a-1)^{-W}\dd p,
\end{equation}
where $c$ satisfies \eqref{rlc:eq:cstrip}. This formula extends it
holomorphically to $|\Ima\mu|<r\pi$.
Every fixed mixed order derivative can be passed through
\eqref{rlc:eq:convolution}; a derivative of total degree $M$ gives
$\partial_W^M\CC_{r,a}(W;\mu)$.
\end{theorem}
\begin{proof}
Put $g_j(x)=e^{-cx}f_{a,s_j}(x)$. The factor $e^{c\mu}$ is constant
on the constraint hyperplane. The integral of absolute values in
\eqref{rlc:eq:convolution} is therefore at most
\[
 e^{c\mu}\|g_r\|_\infty\prod_{j<r}\|g_j\|_1.
\]
Lemma~\ref{rlc:lem:weighted} proves absolute convergence and locally
uniform bounds for all parameter derivatives. The Fourier transform of
$g_j$ is \eqref{rlc:eq:transform} at $p=c+it$. Its product is
\[
 \left(\frac{\pi}{\sin\pi p}\right)^r(p+a-1)^{-W}.
\]
This depends only on $W$. Fourier inversion is applicable because the
product decays like $e^{-r\pi|t|}$ times a fixed power of $1+|t|$ on
compact order sets. It gives \eqref{rlc:eq:barnes} and establishes
pointwise equality, not merely equality almost everywhere; the
convolutions are continuous. The same exponential decay proves the
stated $\mu$ strip and entire dependence on $W$. Norm holomorphy and
Cauchy's formula justify all differentiated integrals.
\end{proof}

For $\Rea W>0$, another useful form follows from the ordinary
one-sided fractional-integral semigroup:
\begin{equation}\label{rlc:eq:basefractional}
 \CC_{r,a}(W;\mu)=\frac1{\Gamma(W)}
 \int_0^\infty u^{W-1}e^{-(a-1)u}
           f_0^{*r}(\mu-u)\dd u.
\end{equation}
One can prove it first with all $\Rea s_j>0$ by Fubini and the
Dirichlet simplex integral; their sum of positive integration variables
has Gamma density of order $W$. Alternatively its transform is
\eqref{rlc:eq:barnes}. The kernel in
\eqref{rlc:eq:basefractional} is then continued through $W=0,-1,\ldots$
by the already established entire function, not by substituting into an
improper integral that no longer converges.

\subsection{A direct two-factor proof}
It is useful to verify the main normalization without Fourier inversion.
For $\Rea s,\Rea t>0$, insert the two Euler integrals and integrate the
reciprocal variable first. For $u,v>0$,
\[
 \int_0^\infty
 \frac{\dd y/y}{(1+ye^{-u})(1+y^{-1}e^{-v})}
 =\frac{u+v}{1-e^{-(u+v)}}.
\]
This follows by one partial fraction decomposition; the logarithms at
infinity cancel. With $w=u+v$, the remaining beta integral yields
\begin{align*}
 \int_0^\infty F_{a,s}(y)F_{a,t}(1/y)\frac{\dd y}{y}
 &=\frac1{\Gamma(W)}\int_0^\infty
       \frac{w^W e^{-(a-1)w}}{e^w-1}\dd w\\
 &=W\zeta(W+1,a).
\end{align*}
Absolute convergence and entire-order continuation proved above give
\eqref{rlc:eq:opening} for all $s,t$. In particular,
\begin{equation}\label{rlc:eq:polylogpair}
 \int_0^\infty\Li_s(-y)\Li_t(-1/y)\frac{\dd y}{y}
     =(s+t)\zeta(s+t+1),
\end{equation}
with the removable value at zero understood. For example, the integrals
for the order pairs $(1,1)$, $(2,1)$, and $(2,2)$ are respectively
$2\zeta(3)$, $3\zeta(4)$, and $4\zeta(5)$.
These examples illustrate the normalization; they are not offered as
historical-priority claims.

\section{Central-factorial kernels and finite zeta reductions}
\label{rlc:sec:central}

\subsection{The logistic convolution polynomial}
Define monic polynomials
\begin{align}
 p_{2m}(X)&=X\prod_{j=1}^{m-1}\bigl(X^2+4j^2\pi^2\bigr),
          &&m\ge1,\label{rlc:eq:peven}\\
 p_{2m+1}(X)&=\prod_{j=1}^{m}\bigl(X^2+(2j-1)^2\pi^2\bigr),
          &&m\ge0.\label{rlc:eq:podd}
\end{align}
Thus $p_1=1$, $p_2=X$, and
\begin{equation}\label{rlc:eq:precurrence}
 p_{r+2}(X)=(X^2+r^2\pi^2)p_r(X).
\end{equation}
These are central-factorial product polynomials after an imaginary
rescaling of their variable. Their degree is $r-1$ and their parity
is $(-1)^{r-1}$.

\begin{lemma}[All-depth elementary kernel]\label{rlc:lem:logistic}
For every $r\ge1$ and real $\mu$,
\begin{equation}\label{rlc:eq:logistic}
 f_0^{*r}(\mu)=
 \frac{p_r(\mu)}{(r-1)!\,[1-(-1)^r e^{-\mu}]}.
\end{equation}
At even depth the quotient at $\mu=0$ is filled in by continuity.
\end{lemma}
\begin{proof}
At depth one this is $f_0$. At depth two the substitution $y=e^x$
gives $f_0*f_0(\mu)=\mu/(1-e^{-\mu})$, including its limit $1$.
Let $g(p)=\pi\csc(\pi p)$. Direct differentiation gives
\[
 \frac{\dd^2}{\dd p^2}g(p)^r
  =r(r+1)g(p)^{r+2}-r^2\pi^2g(p)^r.
\]
Multiplication by $\mu^2$ corresponds to two derivatives of the
bilateral transform. Transform uniqueness therefore implies
\[
 f_0^{*(r+2)}(\mu)=
 \frac{\mu^2+r^2\pi^2}{r(r+1)}f_0^{*r}(\mu).
\]
The two initial cases and \eqref{rlc:eq:precurrence} prove the result.
All transforms used here lie in the common strip $0<\Rea p<1$;
polynomial multiplication is justified by the neighboring-strip
estimates from Lemma~\ref{rlc:lem:weighted}.
\end{proof}

\subsection{A finite formula away from the central point}
Put $\epsilon_r=(-1)^r$.
\begin{theorem}[Finite Lerch closure]\label{rlc:thm:lerch}
For $\mu<0$, $\Rea a>0$, and $W\in\C$,
\begin{equation}\label{rlc:eq:lerchclosure}
 \boxed{\displaystyle
 \CC_{r,a}(W;\mu)=
 \frac{-\epsilon_r e^\mu}{(r-1)!}
 \sum_{j=0}^{r-1}\frac{(-1)^j p_r^{(j)}(\mu)}{j!}
 (W)_j\Phi(\epsilon_r e^\mu,W+j,a).}
\end{equation}
The combined expression continues to all real $\mu$. At even depth and
$\mu>0$, the continuation of the \emph{whole sum} is independent of
whether the Lerch cut is approached from above or below.
\end{theorem}
\begin{proof}
For $\Rea W>0$, insert \eqref{rlc:eq:logistic} into
\eqref{rlc:eq:basefractional}. When $\mu<0$, expand the denominator as
\[
 \frac1{1-\epsilon_r e^{u-\mu}}
 =-\epsilon_r e^{\mu-u}
       \sum_{n=0}^\infty\epsilon_r^n e^{n(\mu-u)}.
\]
Expand $p_r(\mu-u)$ by its finite Taylor formula. The remaining
integrals are
$\Gamma(W+j)(n+a)^{-W-j}/\Gamma(W)$, yielding
\eqref{rlc:eq:lerchclosure}. Absolute convergence justifies the steps.
Both sides are entire in $W$ for fixed $\mu<0$.

The left side is holomorphic for $|\Ima\mu|<r\pi$. At even depth,
the discontinuity of the Lerch integral at $z=e^\mu>1$ is, up to a
common nonzero sign convention,
$2\pi i e^{-a\mu}\mu^{W+j-1}/\Gamma(W+j)$.
After multiplication by $(W)_j$ and the other factors, the jump of the
sum is proportional to
\[
 \sum_{j=0}^{r-1}\frac{(-1)^j p_r^{(j)}(\mu)\mu^j}{j!}
 =p_r(0)=0.
\]
This calculation is initially made away from Gamma poles and then
continued in $W$. Thus the two real boundary values agree and coincide
with the ordinary convolution. At $\mu=0$ use the entire convolution,
or the central formula below. At odd depth the argument
$-e^\mu$ does not cross the positive Lerch cut.
\end{proof}

At depth two and $a=1$, the formula becomes
\begin{equation}\label{rlc:eq:fugacity2}
 \int_0^\infty\Li_s(-y)\Li_t(-e^\mu/y)\frac{\dd y}{y}
 =W\Li_{W+1}(e^\mu)-\mu\Li_W(e^\mu).
\end{equation}
This is a branch-cancelled combination, not a license to discard the
imaginary parts of its terms separately. Its limit at $\mu=0$ is
$\E_1(W)$.

\subsection{Pole-free central formulas}
Define the entire alternating Hurwitz function
\begin{equation}\label{rlc:eq:eta}
 \eta(v,a)=\Phi(-1,v,a)
 =2^{-v}\left[\zeta\left(v,\frac a2\right)
                -\zeta\left(v,\frac{a+1}2\right)\right].
\end{equation}
The two residues at $v=1$ cancel. In particular
$\eta(0,a)=1/2$ and
$\eta(1,a)=\tfrac12[\psi((a+1)/2)-\psi(a/2)]$.
Write
\[
 p_{2m}(X)=\sum_{k=0}^{m-1}A_{m,k}X^{2k+1},\qquad
 p_{2m+1}(X)=\sum_{k=0}^{m}B_{m,k}X^{2k}.
\]
The coefficients are explicitly the elementary symmetric functions
of the squared even or odd multiples of $\pi$ in
\eqref{rlc:eq:peven}--\eqref{rlc:eq:podd}.

\begin{theorem}[Central finite closure]\label{rlc:thm:central}
For $\Rea a>0$ and every $W\in\C$,
\begin{align}
 \CC_{2m,a}(W;0)
 &=\frac1{(2m-1)!}\sum_{k=0}^{m-1}
       A_{m,k}(W)_{2k}\E_a(W+2k),\label{rlc:eq:evencentral}\\
 \CC_{2m+1,a}(W;0)
 &=\frac1{(2m)!}\sum_{k=0}^{m}
       B_{m,k}(W)_{2k}\eta(W+2k,a).\label{rlc:eq:oddcentral}
\end{align}
Every summand in these formulas is entire in $W$.
\end{theorem}
\begin{proof}
For $\Rea W>0$, \eqref{rlc:eq:basefractional} at zero uses
$p_{2m}(u)/(e^u-1)$ in even depth and $p_{2m+1}(u)/(e^u+1)$
in odd depth. The ordinary Mellin integrals are
$(W)_j\zeta(W+j,a)$ or $(W)_j\eta(W+j,a)$.
In even depth $j=2k+1$, so
\[
 (W)_{2k+1}\zeta(W+2k+1,a)
     =(W)_{2k}\E_a(W+2k).
\]
This supplies the pole-free form. Entire continuation proves both
identities at all remaining orders.
\end{proof}

The first cases, useful for checking every normalization, are
\begin{align}
 \CC_{2,a}(W;0)&=\E_a(W),\label{rlc:eq:C2}\\
 \CC_{3,a}(W;0)&=\tfrac12\{(W)_2\eta(W+2,a)+\pi^2\eta(W,a)\},
 \label{rlc:eq:C3}\\
 \CC_{4,a}(W;0)&=\tfrac16\{(W)_2\E_a(W+2)+4\pi^2\E_a(W)\},
 \label{rlc:eq:C4}\\
 \CC_{5,a}(W;0)&=\tfrac1{24}\{(W)_4\eta(W+4,a)
                +10\pi^2(W)_2\eta(W+2,a)+9\pi^4\eta(W,a)\},
 \label{rlc:eq:C5}\\
 \CC_{6,a}(W;0)&=\tfrac1{120}\{(W)_4\E_a(W+4)
                +20\pi^2(W)_2\E_a(W+2)+64\pi^4\E_a(W)\}.
 \label{rlc:eq:C6}
\end{align}
An equivalent differential-operator form is
\begin{equation}\label{rlc:eq:shiftoperator}
 \CC_{r,a}(W;0)=\frac1{(r-1)!}
 \begin{cases}
 p_r(-\partial_a)\zeta(W,a),&r\text{ even},\\
 p_r(-\partial_a)\eta(W,a),&r\text{ odd}.
 \end{cases}
\end{equation}
The even polynomial has no constant term, so it annihilates the
$a$-independent zeta residue.

\subsection{Negative-order resonance is not zero substitution}
At $W=-2$, \eqref{rlc:eq:C4} gives
\begin{equation}\label{rlc:eq:resonance}
 \boxed{\displaystyle
 \CC_{4,a}(-2;0)=\frac13+\frac{2\pi^2}{3}B_2(a),
 \qquad B_2(a)=a^2-a+\frac16.}
\end{equation}
Indeed $(-2)_2\E_a(0)/6=1/3$, while
$\E_a(-2)=-2\zeta(-1,a)=B_2(a)$.
The uncombined expression $(W)_3\zeta(W+3,a)$ has a vanishing
Pochhammer factor next to a pole at $W=-2$. Setting the first factor to
zero loses the $1/3$. This is a concrete integration safeguard, not an
assertion that the inspected repository already printed that error.

\section{Every logarithmic moment at arbitrary complex orders}
\label{rlc:sec:moments}

Order addition alone evaluates the unweighted convolution. Polynomial
weights record how the total logarithmic argument is distributed among
its factors. Surprisingly, they still admit finite closure at arbitrary
complex orders.

\subsection{An all-depth differential-polynomial algorithm}
For $\ell\ge0$ define polynomials $R_\ell(s;T,U)$ by
\begin{align}
 R_0&=1,\notag\\
 R_{\ell+1}&=(T+sU)R_\ell+(\pi^2+T^2)\partial_T R_\ell
                        +U^2\partial_U R_\ell.
 \label{rlc:eq:Rrecurrence}
\end{align}
Their first values are
\begin{align*}
 R_1&=T+sU,\\
 R_2&=2T^2+2sTU+s(s+1)U^2+\pi^2,\\
 R_3&=6T^3+6sT^2U+3s(s+1)TU^2+(s)_3U^3
          +5\pi^2T+3s\pi^2U.
\end{align*}
For a multi-index $\boldsymbol\ell=(\ell_1,\ldots,\ell_r)$, let
\begin{equation}\label{rlc:eq:weighteddef}
 \MM_{\boldsymbol\ell}(\boldsymbol s;a;\mu)
 =\int_{\R^{r-1}}\prod_{j=1}^r
     x_j^{\ell_j}f_{a,s_j}(x_j)\dd x_1\cdots\dd x_{r-1},
 \qquad x_r=\mu-\sum_{j<r}x_j.
\end{equation}

\begin{theorem}[Finite closure of all mixed moments]
\label{rlc:thm:allmoments}
The integral \eqref{rlc:eq:weighteddef} is absolutely convergent and
entire in its orders. Form the finite polynomial
$\prod_{j=1}^r R_{\ell_j}(s_j;T,U)$ and put $W=\sum s_j$.
Apply the following linear replacement rules to its monomials:
\begin{align}
 T^{2k}U^d\ &\longmapsto\
 \sum_{q=0}^k\binom{k}{q}(-\pi^2)^{k-q}
        \CC_{r+2q,a}(W+d;\mu),\label{rlc:eq:evenrule}\\
 T^{2k+1}U^d\ &\longmapsto\
 \sum_{q=0}^k\binom{k}{q}\frac{(-\pi^2)^{k-q}}{r+2q}
 \left\{\mu\CC_{r+2q,a}(W+d;\mu)
 -(W+d)\CC_{r+2q,a}(W+d+1;\mu)\right\}.
 \label{rlc:eq:oddrule}
\end{align}
The resulting expression is exactly
$\MM_{\boldsymbol\ell}(\boldsymbol s;a;\mu)$.
Thus every mixed logarithmic moment at $\mu=0$ reduces to finitely many
entire Hurwitz-zeta combinations in even depth, and alternating
Hurwitz-zeta combinations in odd depth. Arbitrary mixed order
derivatives of this finite formula are valid.
\end{theorem}
\begin{proof}
Put $g(p)=\pi\csc(\pi p)$, $h=p+a-1$, $T=\pi\cot(\pi p)$,
and $U=h^{-1}$. Then
$g'=-Tg$, $T'=-(\pi^2+T^2)$, and $U'=-U^2$.
Repeated differentiation gives the exact identity
\[
 (-\partial_p)^\ell\{g(p)h^{-s}\}
       =g(p)h^{-s}R_\ell(s;T,U).
\]
The left side is the transform of $x^\ell f_{a,s}(x)$.
The product of the transforms of all factors is consequently
$g(p)^r h^{-W}\prod_jR_{\ell_j}(s_j;T,U)$.

The identity $T^2=g^2-\pi^2$ proves
\eqref{rlc:eq:evenrule}. For an odd power, first use that identity
and then $g^R T=-(g^R)'/R$. Integration by parts on the vertical line
in \eqref{rlc:eq:barnes} gives, for every $v$,
\[
 \frac1{2\pi i}\int e^{p\mu}g(p)^R T h^{-v}\dd p
 =\frac{\mu}{R}\CC_{R,a}(v;\mu)
       -\frac{v}{R}\CC_{R,a}(v+1;\mu).
\]
The boundary terms vanish exponentially. This proves the odd rule.
Lemma~\ref{rlc:lem:weighted} supplies absolute convergence and parameter
holomorphy, so no formal interchange remains unproved.
\end{proof}

For example, at $\mu=0$, with $s=s_1$ and $W=\sum s_j$,
\begin{align}
 \MM_{(1,0,\ldots,0)}
 &=\left(s-\frac Wr\right)\CC_{r,a}(W+1;0),
 \label{rlc:eq:firstall}\\
 \MM_{(2,0,\ldots,0)}
 &=2\CC_{r+2,a}(W;0)-\pi^2\CC_{r,a}(W;0)\notag\\
 &\quad+\left[s(s+1)-\frac{2s(W+1)}r\right]
                \CC_{r,a}(W+2;0).
 \label{rlc:eq:secondall}
\end{align}
For two distinct positions having orders $s$ and $t$,
\begin{align}
 \MM_{(1,1,0,\ldots,0)}
 &=\CC_{r+2,a}(W;0)-\pi^2\CC_{r,a}(W;0)\notag\\
 &\quad+\left[st-\frac{(s+t)(W+1)}r\right]
                    \CC_{r,a}(W+2;0).
 \label{rlc:eq:mixedsecond}
\end{align}
Summing \eqref{rlc:eq:firstall} over all positions gives zero, as it
must from $x_1+\cdots+x_r=0$. The factor $1/r$ is essential.

\subsection{A closed two-factor formula}
For $k\ge1$ define
\begin{equation}\label{rlc:eq:Q}
 Q_k(s,t)=\frac{(s)_k-(-1)^k(t)_k}{s+t+k-1}.
\end{equation}
This is a polynomial, not a meromorphic quotient: at
$t=1-k-s$, $(t)_k=(-1)^k(s)_k$, proving divisibility.
Its first values are
\[
 Q_1=1,\qquad Q_2=s-t,\qquad
 Q_3=s^2-st+t^2+s+t.
\]
Let $d_{2j}$ be defined by
\begin{equation}\label{rlc:eq:dcoeff}
 \frac{\pi\lambda}{\sin\pi\lambda}
 =\sum_{j=0}^\infty d_{2j}\lambda^{2j},\qquad
 d_0=1,\quad
 d_{2j}=2(1-2^{1-2j})\zeta(2j)\quad(j\ge1).
\end{equation}

\begin{theorem}[All two-factor logarithmic moments]
\label{rlc:thm:twomoments}
For $\ell\ge0$, $s,t\in\C$, and $\Rea a>0$, put $W=s+t$. Then
\begin{equation}\label{rlc:eq:momentclosed}
 \boxed{\displaystyle
 \int_0^\infty\log^\ell y\,F_{a,s}(y)F_{a,t}(1/y)\frac{\dd y}{y}
 =\ell!\sum_{j=0}^{\lfloor\ell/2\rfloor}
 \frac{d_{2j}}{(\ell-2j+1)!}
 Q_{\ell-2j+1}(s,t)\E_a(W+\ell-2j).}
\end{equation}
Every term on the right is entire in the two independent orders.
\end{theorem}
\begin{proof}
Introduce the locally holomorphic tilted integral
\[
 I(\lambda)=\int_0^\infty y^\lambda
 F_{a,s}(y)F_{a,t}(1/y)\frac{\dd y}{y}.
\]
The weighted estimates prove convergence for
$|\Rea\lambda|<\min(1,\Rea a)$. For
$\Rea s,\Rea t>0$ and sufficiently small $|\lambda|$, inserting
Euler integrals and evaluating the elementary rational $y$ integral
gives
\begin{equation}\label{rlc:eq:tilted}
 I(\lambda)=\frac{\pi}{\sin\pi\lambda}
 \sum_{n=0}^\infty\left[
 (n+a-\lambda)^{-s}(n+a)^{-t}
 -(n+a)^{-s}(n+a+\lambda)^{-t}\right].
\end{equation}
For clarity, the elementary integration kernel before the two Laplace
integrals are performed is
\[
 \frac{\pi}{\sin\pi\lambda}
 \frac{e^{\lambda u}-e^{-\lambda v}}{e^{u+v}-1}
 e^{-(a-1)(u+v)}.
\]
It is multiplied by $u^{s-1}v^{t-1}/(\Gamma(s)\Gamma(t))$.
Expanding the geometric denominator and integrating $u,v$ gives
\eqref{rlc:eq:tilted}. Its $\lambda=0$ value is removable.

Taylor expansion in $\lambda$ now yields
\[
 I(\lambda)=\frac{\pi}{\sin\pi\lambda}
 \sum_{k=1}^\infty
 \frac{(s)_k-(-1)^k(t)_k}{k!}
       \zeta(W+k,a)\lambda^k.
\]
Normal convergence in a small disk follows from $\Rea W>0$ and the
shift bound $|\lambda|<\Rea a$. Extracting the coefficient of
$\lambda^\ell$ and using
\[
 \bigl((s)_k-(-1)^k(t)_k\bigr)\zeta(W+k,a)
       =Q_k(s,t)\E_a(W+k-1)
\]
proves \eqref{rlc:eq:momentclosed} in the initial domain. Both sides
are entire in $s,t$, which proves the theorem everywhere.
\end{proof}

The first two nonconstant moments are therefore
\begin{align}
 \int_0^\infty\log y\,F_{a,s}(y)F_{a,t}(1/y)\frac{\dd y}{y}
 &=\frac{s-t}{2}\E_a(W+1),\label{rlc:eq:moment1}\\
 \int_0^\infty\log^2 y\,F_{a,s}(y)F_{a,t}(1/y)\frac{\dd y}{y}
 &=\frac{s^2-st+t^2+s+t}{3}\E_a(W+2)
      +\frac{\pi^2}{3}\E_a(W).
 \label{rlc:eq:moment2}
\end{align}
In particular, at $s=t=0$ the second moment is $\pi^2/3$, the ordinary
logistic second moment. Formula \eqref{rlc:eq:moment1} independently
checks the $1/2$ in the first-moment normalization. All odd moments
vanish when $s=t$, directly reflecting $y\leftrightarrow1/y$.

\section{Ordinary Stieltjes products and harmonic coefficient calculus}
\label{rlc:sec:stieltjes}

\subsection{A convergent product at every Stieltjes index}
Set
\begin{equation}\label{rlc:eq:Ljets}
 L_m^a(y)=\left.\partial_s^mF_{a,s}(y)\right|_{s=0}.
\end{equation}
These are order derivatives, not argument derivatives. From
\eqref{rlc:eq:stieltjesconvention},
\begin{equation}\label{rlc:eq:Ejets}
 \E_a^{(M)}(0)=
 \begin{cases}
 1,&M=0,\\
 (-1)^{M-1}M\gamma_{M-1}(a),&M\ge1.
 \end{cases}
\end{equation}

\begin{theorem}[Ordinary Stieltjes product law]
\label{rlc:thm:stieltjes}
For $m,n\ge0$ and $\Rea a>0$,
\begin{equation}\label{rlc:eq:stieltjespair}
 \boxed{\displaystyle
 \int_0^\infty L_m^a(y)L_n^a(1/y)\frac{\dd y}{y}
 =\begin{cases}
 1,&m=n=0,\\
 (-1)^{m+n-1}(m+n)\gamma_{m+n-1}(a),&m+n>0.
 \end{cases}}
\end{equation}
All integrals are absolutely convergent.
\end{theorem}
\begin{proof}
Differentiate \eqref{rlc:eq:opening} $m$ times in $s$ and $n$ times
in $t$. Lemma~\ref{rlc:lem:weighted} justifies the differentiation.
The right side is \eqref{rlc:eq:Ejets} with $M=m+n$.
\end{proof}

For example,
\[
 \int_0^\infty L_1^a(y)L_0^a(1/y)\frac{\dd y}{y}
 =-\psi(a),\qquad
 \int_0^\infty L_1^a(y)L_1^a(1/y)\frac{\dd y}{y}
 =-2\gamma_1(a).
\]
At $a=1$ these become identities involving only order derivatives
of ordinary polylogarithms. The apparent product can have either sign:
this is not a Hermitian inner product and the second factor has
reciprocal argument rather than complex conjugation.

\subsection{Logarithmic Stieltjes identities}
Theorem~\ref{rlc:thm:twomoments} can be differentiated independently in
its two orders. Three explicit consequences are
\begin{align}
 \int_0^\infty\log y\,L_1^a(y)L_0^a(1/y)\frac{\dd y}{y}
 &=\frac12\zeta(2,a),\label{rlc:eq:jetlog1}\\
 \int_0^\infty\log^2 y\,L_1^a(y)L_0^a(1/y)\frac{\dd y}{y}
 &=\frac{\pi^2}{3}\gamma_0(a)+\frac23\zeta(3,a),
 \label{rlc:eq:jetlog2}\\
 \int_0^\infty\log^2 y\,L_1^a(y)L_1^a(1/y)\frac{\dd y}{y}
 &=\frac43\zeta^{[1]}(3,a)-\frac{2\pi^2}{3}\gamma_1(a).
 \label{rlc:eq:jetlogmixed}
\end{align}
To check the last identity explicitly, write
$Q=s^2-st+t^2+s+t$ and differentiate
$Q\E_a(s+t+2)$ at $s=t=0$. The result is
$-\E_a(2)+2\E_a'(2)=4\zeta^{[1]}(3,a)$.
The remaining term is $\pi^2\E_a''(0)/3$.
This derivation fixes both the coefficient $4/3$ and the Stieltjes sign.

\subsection{All-depth jets as finite harmonic polynomials}
For a multi-index $\boldsymbol m$ of length $r$, let
$J_{r,\boldsymbol m}(a)$ be the integral at $\mu=0$ of the product of
$f_{a,s_j}^{[m_j]}(x_j)$ evaluated at all $s_j=0$. Let
$M=\sum m_j$. Order addition proves that this integral depends only on
$M$, so write it as $J_{r,M}(a)$.
The next result makes that value completely explicit.

\begin{theorem}[All-depth Stieltjes coefficients]
\label{rlc:thm:alljets}
At even depth, for $M\ge1$,
\begin{align}
 J_{2m,M}(a)=\frac1{(2m-1)!}\biggl[&
 A_{m,0}(-1)^{M-1}M\gamma_{M-1}(a)\notag\\
 &+\sum_{k=1}^{m-1}A_{m,k}
 \sum_{h=1}^{\min(M,2k+1)}
 \frac{M!}{(M-h)!}\stir{2k+1}{h}
       \zeta^{[M-h]}(2k+1,a)\biggr].
 \label{rlc:eq:evenjets}
\end{align}
Here $\stir{j}{h}$ is an unsigned Stirling number of the first kind.
For $M=0$, $J_{2m,0}=A_{m,0}/(2m-1)!$.
At odd depth the corresponding formula is
\begin{equation}\label{rlc:eq:oddjets}
 J_{2m+1,M}(a)=\frac1{(2m)!}
 \sum_{k=0}^{m}B_{m,k}
 \sum_{h=0}^{\min(M,2k)}
 \frac{M!}{(M-h)!}\stir{2k}{h}
          \eta^{[M-h]}(2k,a),
\end{equation}
where $\stir00=1$. In particular $J_{2m+1,0}=B_{m,0}/[2(2m)!]$.
\end{theorem}
\begin{proof}
Differentiate \eqref{rlc:eq:evencentral} and
\eqref{rlc:eq:oddcentral} $M$ times at zero. In even depth keep the
$k=0$ term in the entire form $A_{m,0}\E_a(W)$. For $k\ge1$ it is
safe to use $(W)_{2k+1}\zeta(W+2k+1,a)$.
The identity $(W)_j=\sum_h\stir jh W^h$ and the ordinary product rule
give \eqref{rlc:eq:evenjets}. The same expansion gives
\eqref{rlc:eq:oddjets}. The zero-order formulas follow immediately.
\end{proof}

For example,
\begin{align}
 J_{4,1}(a)&=\frac{2\pi^2}{3}\gamma_0(a)+\frac13\zeta(3,a),
 \label{rlc:eq:J41}\\
 J_{4,2}(a)&=-\frac{4\pi^2}{3}\gamma_1(a)
              +\zeta(3,a)+\frac23\zeta^{[1]}(3,a).
 \label{rlc:eq:J42}
\end{align}
The number of factors and the distribution of the derivative orders
can be varied without changing these scalar values, provided their
respective totals stay fixed.

The coefficients also have a finite generalized-harmonic description.
For $j\ge1$,
\begin{equation}\label{rlc:eq:pochhammerharmonic}
 (z)_j=(j-1)!z\exp\left(
 \sum_{q=1}^\infty\frac{(-1)^{q-1}}q H_{j-1}^{(q)}z^q\right),
\end{equation}
locally at zero. Thus, for $1\le h\le j$,
\begin{equation}\label{rlc:eq:Bellharmonic}
 \left.\partial_z^h(z)_j\right|_{0}
 =h(j-1)!\,\mathcal B_{h-1}
 \left(H_{j-1},-H_{j-1}^{(2)},2!H_{j-1}^{(3)},\ldots\right),
\end{equation}
where $\mathcal B_n$ is the complete exponential Bell polynomial.
Indeed, divide $(z)_j$ by $(j-1)!z$, take the logarithm of
$\prod_{q=1}^{j-1}(1+z/q)$, and expand. This proves the formula,
including the signs and factorials. At nonzero resonances one should
use the finite polynomial $(z)_j$ itself, rather than a logarithmic
derivative divided by a vanishing factor.

Together with Theorem~\ref{rlc:thm:allmoments}, these formulas give a
terminating algorithm for every mixed order jet of every logarithmic
moment. The coefficient operations are finite; they do not presuppose
arithmetic independence of the resulting zeta or Stieltjes values.

\section{Reciprocal remainders and Gaussian odd parts}
\label{rlc:sec:harmonic}

\subsection{An all-order identity for polylogarithmic remainders}
For $N\ge0$ define the finite Taylor remainder
\begin{equation}\label{rlc:eq:remainderdef}
 \mathcal R_{N,s}(y)=\Li_s(-y)-\sum_{k=1}^{N}\frac{(-y)^k}{k^s}.
\end{equation}
The empty sum at $N=0$ is zero. The elementary Lerch shift formula
implies
\begin{equation}\label{rlc:eq:remaindershift}
 F_{N+1,s}(y)=-(-y)^{-N}\mathcal R_{N,s}(y).
\end{equation}
Although each remainder can grow polynomially on the positive real
axis, the reciprocal product is integrable.

\begin{corollary}[Harmonic-tail product identity]
\label{rlc:cor:remainders}
For every $N\ge0$ and $s,t\in\C$,
\begin{equation}\label{rlc:eq:remainderidentity}
 \boxed{\displaystyle
 \int_0^\infty\mathcal R_{N,s}(y)\mathcal R_{N,t}(1/y)
                  \frac{\dd y}{y}
 =\E_{N+1}(s+t)
 =(s+t)\left[\zeta(s+t+1)-H_N^{(s+t+1)}\right].}
\end{equation}
The last expression is filled in by its removable value $1$ when
$s+t=0$. Every order derivative and every logarithmic moment is obtained
by substituting $a=N+1$ into the preceding theorems.
\end{corollary}
\begin{proof}
The two powers $(-y)^{-N}$ and $(-1/y)^{-N}$ in
\eqref{rlc:eq:remaindershift} multiply to one. Apply
\eqref{rlc:eq:opening}, followed by the Hurwitz shift identity
$\zeta(v,N+1)=\zeta(v)-H_N^{(v)}$. Normal convergence already proved
for the Lerch profiles gives all differentiated versions.
\end{proof}

At $N=1$ and $s=t=1$, this gives the elementary-looking but nontrivial
identity
\begin{equation}\label{rlc:eq:remainderexample}
 \int_0^\infty
 [\log(1+y)-y][\log(1+1/y)-1/y]\frac{\dd y}{y}
       =2\zeta(3)-2.
\end{equation}
The factors must remain combined. Expanding their product and
integrating the resulting pieces independently introduces divergent
integrals and is not a valid proof of the formula.

The Stieltjes constants at the shifted integer arguments also reduce
by a finite harmonic-logarithmic correction:
\begin{equation}\label{rlc:eq:stieltjesshift}
 \gamma_n(N+1)=\gamma_n(1)
           -\sum_{k=1}^N\frac{\log^n k}{k}.
\end{equation}
It follows by comparing Laurent coefficients of
$\zeta(1+v,N+1)=\zeta(1+v)-\sum_{k=1}^N k^{-1-v}$.
For $n=0$ it reads $\gamma_0(N+1)=\gamma-H_N$.
Thus the entire remainder family contains finite generalized harmonic
numbers and their order derivatives without changing the convergence
normalization.

\subsection{A fourth-root polylogarithm family}
For complex $s$ and positive $y$, define the analytic odd part
\begin{equation}\label{rlc:eq:Adef}
 A_s(y)=\frac{\Li_s(iy)-\Li_s(-iy)}{2i}.
\end{equation}
For real $s$ this agrees with an imaginary part. For complex $s$,
\eqref{rlc:eq:Adef}, rather than $\Ima\Li_s(iy)$, is required for
holomorphic order differentiation.
The initial series gives
\[
 A_s(y)=\sum_{n=0}^\infty\frac{(-1)^ny^{2n+1}}{(2n+1)^s},\qquad
 F_{1/2,s}(y^2)=2^s y A_s(y),
\]
with continuation to the positive real axis.

\begin{corollary}[Gaussian reciprocal order addition]
\label{rlc:cor:gaussian}
For all $s,t\in\C$, with $W=s+t$,
\begin{equation}\label{rlc:eq:gaussian}
 \boxed{\displaystyle
 \int_0^\infty A_s(y)A_t(1/y)\frac{\dd y}{y}
 =2^{-W-1}\E_{1/2}(W)
       =(1-2^{-W-1})\E_1(W).}
\end{equation}
Every mixed order derivative is an ordinary convergent integral.
\end{corollary}
\begin{proof}
Use the displayed relation to $F_{1/2,s}$ and set $z=y^2$.
The measure contributes $1/2$ and the order prefactors contribute
$2^{-W}$. Theorem~\ref{rlc:thm:addition} at depth two gives the first
equality. The second follows from
$\zeta(v,1/2)=(2^v-1)\zeta(v)$, with the pole cancelled before evaluation.
\end{proof}

Since $A_1(y)=\arctan y$, a basic specialization is
\begin{equation}\label{rlc:eq:arctan}
 \int_0^\infty\arctan y\,\arctan(1/y)\frac{\dd y}{y}
     =\frac74\zeta(3).
\end{equation}
More characteristic of the order-derivative continuation is the identity
\begin{equation}\label{rlc:eq:gaussianjets}
 \int_0^\infty A_0^{[1]}(y)A_0^{[1]}(1/y)\frac{\dd y}{y}
 =-\gamma_1+\gamma\log2-\frac12\log^22.
\end{equation}
Here $A_0^{[1]}=\left.\partial_sA_s\right|_{s=0}$.
It follows by taking the second derivative at zero of
$(1-2^{-W-1})\E_1(W)$.

These identities use fourth-root arguments, but their relation to the
repository's Gaussian depth-two periods is not automatic. In particular,
\eqref{rlc:eq:gaussian} is not a reduction of
$\sum_{n\ge0}(-1)^nH_n/(2n+1)^6$ in the prescribed $g_{a,b}$ basket.
That distinction is retained in the integration status.

\section{Gamma-weighted integrals and explicit polygamma bridges}
\label{rlc:sec:gamma}

When $\Rea a>1/2$, the symmetric vertical line $p=1/2+it$ in
\eqref{rlc:eq:barnes} is admissible. Put $b=a-1/2$. Euler reflection
then gives
\begin{equation}\label{rlc:eq:gammacentral}
 \boxed{\displaystyle
 \begin{aligned}
 \CC_{r,a}(W;0)&=\frac1{2\pi}\int_{-\infty}^{\infty}
 \left[\Gamma\left(\tfrac12+it\right)
       \Gamma\left(\tfrac12-it\right)\right]^r
 (b+it)^{-W}\dd t\\
 &=\frac1{2\pi}\int_\R
      \left(\frac\pi{\cosh\pi t}\right)^r(b+it)^{-W}\dd t.
 \end{aligned}}
\end{equation}
The logarithm of $b+it$ is the right-half-plane branch. For
$0<\Rea a\le1/2$ use a nonsymmetric line satisfying
\eqref{rlc:eq:cstrip}; simply keeping the symmetric line can cross
$\Rea(p+a-1)=0$ and is not the stated formula.

Theorem~\ref{rlc:thm:central} now evaluates every integral in
\eqref{rlc:eq:gammacentral} by a finite Hurwitz-zeta expression.
For $r=2$ and $M\ge1$, differentiation in $W$ at zero gives
\begin{equation}\label{rlc:eq:gammastieltjes}
 \boxed{\displaystyle
 \frac\pi2\int_\R
     \frac{\log^M(a-\tfrac12+it)}{\cosh^2\pi t}\dd t
       =-M\gamma_{M-1}(a).}
\end{equation}
There is no alternating sign on the right: the $(-1)^M$ from the
power derivative cancels all but a minus sign in
\eqref{rlc:eq:Ejets}. This is a consistency bridge to classical
Stieltjes integral representations, not a claim of first discovery
of such representations.

\subsection{An ordinary log-Gamma identity}
At $W=-1$, differentiation of
$\E_a(W)=W\zeta(W+1,a)$ yields
\[
 \E_a'(-1)=\zeta(0,a)-\zeta^{[1]}(0,a).
\]
Using \eqref{rlc:eq:gammacentral} and
$\zeta^{[1]}(0,a)=\log\Gamma(a)-\tfrac12\log(2\pi)$ gives
\begin{equation}\label{rlc:eq:gammalog}
 \frac\pi2\int_\R\frac{
 (b+it)\log(b+it)-(b+it)}{\cosh^2\pi t}\dd t
      =\log\Gamma(a)-\frac12\log(2\pi).
\end{equation}
All differentiations are dominated by the hyperbolic exponential
weight. The classical Hurwitz--Gamma identity is the only additional
input \cite{DLMF}.

\subsection{Higher-depth Gamma and polygamma identities}
The all-depth jets in Section~\ref{rlc:sec:stieltjes} give, for example,
\begin{align}
 \frac1{2\pi}\int_\R
 \left(\frac\pi{\cosh\pi t}\right)^4\log(b+it)\dd t
 &=\frac{2\pi^2}{3}\psi(a)+\frac16\psi^{(2)}(a),
 \label{rlc:eq:gamma4log}\\
 \frac1{2\pi}\int_\R
 \left(\frac\pi{\cosh\pi t}\right)^4\log^2(b+it)\dd t
 &=-\frac{4\pi^2}{3}\gamma_1(a)
     +\zeta(3,a)+\frac23\zeta^{[1]}(3,a).
 \label{rlc:eq:gamma4log2}
\end{align}
The first uses $\psi^{(2)}(a)=-2\zeta(3,a)$.
Positive integer total orders likewise reduce to polygamma values by
\[
 \zeta(k,a)=\frac{(-1)^k}{(k-1)!}\psi^{(k-1)}(a),\qquad k\ge2.
\]
At integer shifts these combine with
$\zeta(k,N+1)=\zeta(k)-H_N^{(k)}$; at rational shifts they give finite
root-of-unity coordinates after the usual residue-class filter.

One can also evaluate every polynomial $t$ moment in
\eqref{rlc:eq:gammacentral}. Expand
$it=(b+it)-b$ to obtain
\begin{equation}\label{rlc:eq:tGamma}
 \frac1{2\pi}\int_\R
 \left(\frac\pi{\cosh\pi t}\right)^r(it)^q(b+it)^{-W}\dd t
 =\sum_{j=0}^{q}\binom qj(-b)^{q-j}\CC_{r,a}(W-j;0).
\end{equation}
Subsequent order derivatives insert arbitrary powers of $\log(b+it)$.
Thus the same finite central-factorial family controls Gamma weights,
polynomial moments, polygamma values, and Stieltjes order derivatives.

\section{Unequal shifts: finite reductions and arithmetic collapses}
\label{rlc:sec:unequal}

Different shifts replace the single resolvent in
\eqref{rlc:eq:barnes} by a product. For $\Rea a_j>0$ define
\[
 \mathcal K(\boldsymbol s,\boldsymbol a;\mu)
       =(f_{a_1,s_1}*\cdots*f_{a_r,s_r})(\mu).
\]
A common $c$ with
$\max_j(0,1-\Rea a_j)<c<1$ exists. The weighted proof gives
\begin{equation}\label{rlc:eq:unequaltransform}
 \mathcal K=\frac1{2\pi i}\int_{c-i\infty}^{c+i\infty}
       e^{p\mu}g(p)^r\prod_{j=1}^r(p+a_j-1)^{-s_j}\dd p,
 \qquad g(p)=\pi\csc(\pi p).
\end{equation}
This is an entire function of the independent orders.

\subsection{A precise rigidity statement}
\begin{proposition}[Rigidity of pure order addition]
\label{rlc:prop:rigidity}
Group equal shifts into distinct values $a_\nu$, with grouped orders
$S_\nu$. Put $W=\sum_\nu S_\nu$. For a proposed shift $b$ with
$\Rea b>0$, the identity
\[
 \mathcal K(\boldsymbol s,\boldsymbol a;\mu)=\CC_{r,b}(W;\mu)
\]
holds on an open interval of real $\mu$ if and only if every grouped
order at a shift different from $b$ is zero and the grouped order at
$b$ is $W$; absent groups are interpreted as zero. In particular, for
strictly positive real orders it forces all shifts to equal $b$.
\end{proposition}
\begin{proof}
Analyticity in the $\mu$ strip turns interval equality into equality on
the whole real axis. Taking the bilateral transform, cancelling $g^r$,
and logarithmically differentiating gives the rational identity
\[
 \sum_\nu\frac{S_\nu}{p+a_\nu-1}=\frac{W}{p+b-1}.
\]
Equality of residues at distinct poles gives the necessity. Conversely,
the stated grouped cancellations give equality of the multipliers with
the chosen branches and hence equality of convolutions. Equivalently,
a possible multiplicative constant after logarithmic differentiation is
one by the behavior as real $p\to+\infty$ of the resolvent product.
\end{proof}

This is a functional rigidity statement within a specified family.
It is not an arithmetic-independence theorem or a prohibition on other
special-value reductions. Broader finite-jet nonclosure by monodromy
has already been studied for the periodic unequal-twist family
\cite{UTJ}; that preceding result is not claimed anew here.

\subsection{Dirichlet mixtures and integer-order partial fractions}
For $\Rea s_j>0$, the classical Dirichlet parameter formula gives
\begin{equation}\label{rlc:eq:simplex}
 \mathcal K(\boldsymbol s,\boldsymbol a;\mu)
 =\frac{\Gamma(W)}{\prod_j\Gamma(s_j)}
 \int_{\substack{u_j\ge0\\\sum u_j=1}}
       \prod_j u_j^{s_j-1}
       \CC_{r,\sum u_ja_j}(W;\mu)\dd u_1\cdots\dd u_{r-1}.
\end{equation}
Indeed apply the Gamma integral to each resolvent in
\eqref{rlc:eq:unequaltransform}, separate the total integration variable
from its simplex proportions, and then invert the transform. The real
parts of the convexly combined shifts remain positive, so Fubini is
justified by a common weighted bound.

If the grouped orders are positive integers $m_\nu$, the mixture has a
finite alternative. For distinct $a_1,\ldots,a_d$ set
\begin{equation}\label{rlc:eq:partialcoeff}
 c_{\nu,k}=\frac1{(m_\nu-k)!}
 \left.\frac{\dd^{m_\nu-k}}{\dd z^{m_\nu-k}}
 \prod_{\rho\ne\nu}(z+a_\rho-a_\nu)^{-m_\rho}\right|_{z=0}.
\end{equation}
The elementary partial fraction identity and transform uniqueness give
\begin{equation}\label{rlc:eq:partialclosure}
 \boxed{\displaystyle
 \mathcal K(\boldsymbol m,\boldsymbol a;\mu)
       =\sum_{\nu=1}^{d}\sum_{k=1}^{m_\nu}
             c_{\nu,k}\CC_{r,a_\nu}(k;\mu).}
\end{equation}
The number $r$ is the number of original profile factors, not the number
of distinct shifts after grouping. At zero, the right side is finite
polygamma or alternating-polygamma data. Colliding shifts are handled
by the limit of the entire combined expression or by grouping first;
individual partial fraction coefficients need not have finite limits.

For instance,
\begin{equation}\label{rlc:eq:unequal11}
 \int_0^\infty F_{a,1}(y)F_{b,1}(1/y)\frac{\dd y}{y}
       =\frac{\zeta(2,a)-\zeta(2,b)}{b-a}.
\end{equation}
At $b=a$ the value is $2\zeta(3,a)$. For unequal nonnegative integers
$N,M$, this specializes to the finite harmonic quotient
\begin{equation}\label{rlc:eq:unequalharmonic}
 \int_0^\infty F_{N+1,1}(y)F_{M+1,1}(1/y)\frac{\dd y}{y}
       =\frac{H_M^{(2)}-H_N^{(2)}}{M-N}.
\end{equation}
In particular the values for $(N,M)=(0,1)$ and $(0,2)$ are $1$ and
$5/8$, respectively.

\subsection{All-depth elementary integrals with finite harmonic values}
Unit-order profiles at integer shifts are elementary logarithmic
remainders by \eqref{rlc:eq:remaindershift}. Their multivariable
convolutions can have especially simple arithmetic values.
For distinct nonnegative integers $N_1,\ldots,N_r$, put
\[
 c_\nu=\prod_{\rho\ne\nu}(N_\rho-N_\nu)^{-1},\qquad
 \overline H_N^{(q)}=\sum_{j=1}^N\frac{(-1)^{j-1}}{j^q},
\]
and let $\mathcal K_{\boldsymbol N}$ denote
$(f_{N_1+1,1}*\cdots*f_{N_r+1,1})(0)$.

\begin{theorem}[Integral-shift arithmetic collapse]
\label{rlc:thm:lattice}
For $r=2m\ge2$,
\begin{equation}\label{rlc:eq:latticeeven}
 \mathcal K_{\boldsymbol N}
 =-\frac1{(2m-1)!}\sum_{k=0}^{m-1}A_{m,k}(2k+1)!
          \sum_{\nu=1}^{2m}c_\nu H_{N_\nu}^{(2k+2)}.
\end{equation}
It belongs to $\mathbb Q[\pi^2]$, of degree at most $m-1$ in $\pi^2$.
For $r=2m+1\ge3$, if all $N_\nu$ have the same parity and
$\sigma=(-1)^{N_1}$, then
\begin{equation}\label{rlc:eq:latticeodd}
 \mathcal K_{\boldsymbol N}
 =-\frac{\sigma}{(2m)!}\sum_{k=0}^{m}B_{m,k}(2k)!
          \sum_{\nu=1}^{2m+1}c_\nu\overline H_{N_\nu}^{(2k+1)}.
\end{equation}
This belongs to $\mathbb Q[\pi^2]$, of degree at most $m$.
\end{theorem}
\begin{proof}
The unit-order partial fraction formula is
$\mathcal K_{\boldsymbol N}=\sum c_\nu\CC_{r,N_\nu+1}(1;0)$.
Moreover $\sum c_\nu=0$ for $r\ge2$, by comparing coefficients of
$p^{-1}$ at infinity in the resolvent identity.

At even depth insert the central zeta formula and
$\zeta(q,N+1)=\zeta(q)-H_N^{(q)}$. The common zeta terms cancel because
$\sum c_\nu=0$, giving \eqref{rlc:eq:latticeeven}.
At odd depth use
\[
 \eta(q,N+1)=(-1)^N\{\eta(q,1)-\overline H_N^{(q)}\}.
\]
The common parity makes the sign independent of $\nu$, so the common
alternating-zeta terms cancel as well. This proves
\eqref{rlc:eq:latticeodd}. All remaining harmonic numbers are rational,
and the degrees in $\pi^2$ follow from the product polynomials.
\end{proof}

Two explicit instances are
\begin{align}
 (f_{1,1}*f_{3,1}*f_{5,1})(0)
       &=\frac{1475}{13824}+\frac{5\pi^2}{192},
       \label{rlc:eq:lattice3}\\
 (f_{1,1}*f_{2,1}*f_{3,1}*f_{4,1})(0)
       &=\frac{575}{3888}+\frac{11\pi^2}{162}.
       \label{rlc:eq:lattice4}
\end{align}
These are, respectively, ordinary two- and three-dimensional integrals
of explicitly elementary logarithmic profiles. Their derivation requires
no integer-relation search.

\subsection{A closed consecutive-shift harmonic formula}
Define complete homogeneous harmonic polynomials $h_j^{(n)}$ by
\begin{equation}\label{rlc:eq:completeharmonic}
 \prod_{q=1}^n(1-z/q)^{-1}=\sum_{j=0}^\infty h_j^{(n)}z^j,
 \qquad
 h_j^{(n)}=\frac1{j!}\mathcal B_j
       \bigl(H_n,1!H_n^{(2)},2!H_n^{(3)},\ldots\bigr).
\end{equation}
These have positive harmonic signs, in contrast to the elementary
product in \eqref{rlc:eq:pochhammerharmonic}.

\begin{corollary}[Consecutive even-depth evaluation]
\label{rlc:cor:consecutive}
For $r=2m$ and $n=2m-1$,
\begin{equation}\label{rlc:eq:consecutive}
 \boxed{\displaystyle
 (f_{1,1}*f_{2,1}*\cdots*f_{2m,1})(0)
   =\frac1{n(n!)^2}\sum_{k=0}^{m-1}
          A_{m,k}(2k+1)!h_{2k+1}^{(n)}.}
\end{equation}
\end{corollary}
\begin{proof}
In \eqref{rlc:eq:latticeeven}, the shifts are $N_\nu=0,1,\ldots,n$
and $c_N=(-1)^N/[N!(n-N)!]$.
Interchanging two finite sums gives
\[
 \sum_{N=0}^n c_NH_N^{(q)}
 =\frac1{n!}\sum_{j=1}^n\frac{(-1)^j}{j^q}\binom{n-1}{j-1}
 =-\frac{h_{q-1}^{(n)}}{n\,n!}.
\]
For the last step use
$\binom{n-1}{j-1}=j\binom nj/n$ and
\[
 \sum_{j=1}^n(-1)^{j-1}\binom nj j^{-v}=h_v^{(n)}.
\]
The latter identity follows by the partial fraction decomposition of
$\prod_{q=1}^n(1-z/q)^{-1}$ and comparison of its $z^v$ coefficient.
Substitution proves \eqref{rlc:eq:consecutive}. The Bell expression in
\eqref{rlc:eq:completeharmonic} follows by taking a logarithm of its
finite product.
\end{proof}

Thus the reciprocal convolution calculus naturally contains both
harmonic coefficient systems: the elementary product governing order
derivatives of rising factorials, and the reciprocal product governing
these finite harmonic collapses. Neither system needs to be inferred
from a numerical pattern.

\section{Anchored antiderivatives and a finite derivative--primitive algebra}
\label{rlc:sec:primitives}

\subsection{The covariant fugacity ladder}
The transform immediately relates changing total order to integrating
or differentiating the logarithmic product argument.

\begin{theorem}[Covariant primitives and shift derivatives]
\label{rlc:thm:ladder}
For $\Rea a>0$, $W\in\C$, and real $\mu$,
\begin{align}
 (\partial_\mu+a-1)\CC_{r,a}(W;\mu)
      &=\CC_{r,a}(W-1;\mu),\label{rlc:eq:muladder}\\
 \CC_{r,a}(W+1;\mu)
      &=e^{-(a-1)\mu}\int_{-\infty}^{\mu}
              e^{(a-1)v}\CC_{r,a}(W;v)\dd v,
        \label{rlc:eq:muprimitive}\\
 \partial_a^n\CC_{r,a}(W;\mu)
      &=(-1)^n(W)_n\CC_{r,a}(W+n;\mu),\qquad n\ge0.
       \label{rlc:eq:aderivative}
\end{align}
The primitive in \eqref{rlc:eq:muprimitive} is fixed by the condition
$e^{(a-1)\mu}\CC_{r,a}(W+1;\mu)\to0$ as $\mu\to-\infty$.
All identities can be differentiated in $W$.
\end{theorem}
\begin{proof}
In \eqref{rlc:eq:barnes}, $\partial_\mu+a-1$ multiplies the integrand
by $p+a-1$, proving \eqref{rlc:eq:muladder}. Differentiating the power
$(p+a-1)^{-W}$ in $a$ gives \eqref{rlc:eq:aderivative}.
The contour bound for real negative $\mu$ is $O(e^{c\mu})$ on compact
order sets. Since $c+\Rea a-1>0$, multiplication by
$e^{(a-1)\mu}$ tends to zero at $-\infty$ and the integral in
\eqref{rlc:eq:muprimitive} converges. Integrating
\eqref{rlc:eq:muladder} at order $W+1$ proves the anchored formula.
The same bounds with logarithmic factors justify all order derivatives.
\end{proof}

At $a=1$, this is an ordinary primitive ladder in $\mu$, or an Euler
primitive ladder in $q=e^\mu$ with measure $\dd q/q$. Raising $W$ by
any positive integer supplies repeated primitives with all constants
fixed. At other shifts the integrating factor is indispensable.

Let $D_a=\partial_\mu+a-1$ and let $J_a$ denote the anchored integral
operator on the right of \eqref{rlc:eq:muprimitive}. On the linear span
of the convolution profiles and their order derivatives,
\begin{equation}\label{rlc:eq:commutators}
 [\partial_a,D_a]=1,\qquad
 [\partial_a,J_a]=-J_a^2.
\end{equation}
The first identity is immediate. For the second, differentiation of the
integrating factor gives
\[
 (\partial_aJ_a-J_a\partial_a)h(\mu)
   =-\int_{-\infty}^{\mu}(\mu-v)e^{-(a-1)(\mu-v)}h(v)\dd v
   =-J_a^2h(\mu).
\]
Absolute convergence justifies the repeated integral. Thus a derivative
in the shift does not commute with a covariant primitive; the correction
is an explicitly normalized second primitive.

Combining \eqref{rlc:eq:aderivative} with order differentiation gives
\begin{equation}\label{rlc:eq:jetderivative}
 \partial_a^n\partial_W^M\CC_{r,a}(W;\mu)
 =(-1)^n\sum_{h=0}^{\min(M,n)}\binom Mh
       \partial_W^h(W)_n\,
       \partial_W^{M-h}\CC_{r,a}(W+n;\mu).
\end{equation}
The coefficients at $W=0$ are precisely the finite harmonic
polynomials in \eqref{rlc:eq:Bellharmonic}. Equations
\eqref{rlc:eq:commutators}--\eqref{rlc:eq:jetderivative}, together with
the reciprocal harmonic product in \eqref{rlc:eq:completeharmonic},
provide a concrete compatible derivative--primitive algebra in this
family. This addresses a specified instance of the compatibility
question in the Mellin continuation \cite{MLD}, not an assertion that
every harmonic calculus in the repository has thereby been unified.

\subsection{Anchored integration in the Hurwitz shift}
For $\Rea a,\Rea b>0$, integration can be taken along any path in the
right half-plane. A pole-free fundamental identity is
\begin{equation}\label{rlc:eq:Eprimitive}
 \int_a^b\E_c(W)\dd c=\zeta(W,a)-\zeta(W,b).
\end{equation}
It follows from $\partial_c\zeta(W,c)=-W\zeta(W+1,c)$.
The right side is entire in $W$, since the residues at $W=1$ cancel.
The central even-depth analogue is
\begin{equation}\label{rlc:eq:evenprimitive}
 \int_a^b\CC_{2m,c}(W;0)\dd c
 =\frac1{(2m-1)!}\sum_{k=0}^{m-1}A_{m,k}(W)_{2k}
      [\zeta(W+2k,a)-\zeta(W+2k,b)].
\end{equation}
Every bracket is taken as an entire combined difference before
specialization.

At odd depth, put $\Delta\eta(v)=\eta(v,a)-\eta(v,b)$. Then
\begin{align}
 \int_a^b\CC_{2m+1,c}(W;0)\dd c
 =\frac1{(2m)!}\biggl[&B_{m,0}\frac{\Delta\eta(W-1)}{W-1}\notag\\
 &+\sum_{k=1}^{m}B_{m,k}(W)_{2k-1}
                          \Delta\eta(W+2k-1)\biggr].
 \label{rlc:eq:oddprimitive}
\end{align}
The first quotient is removable at $W=1$, because $\eta(0,a)=\eta(0,b)$.
Its value there is $\eta^{[1]}(0,a)-\eta^{[1]}(0,b)$.
These two formulas follow by differentiating their right sides and using
the anchor $b=a$; they therefore specify the integration constants as
well as the derivatives.

\subsection{Stieltjes antiderivatives and Gamma normalization}
Taking order derivatives of \eqref{rlc:eq:Eprimitive} at zero gives
\begin{equation}\label{rlc:eq:gammaprimitive}
 \int_a^b\gamma_0(c)\dd c=\log\Gamma(a)-\log\Gamma(b),
\end{equation}
and, for every $n\ge0$,
\begin{equation}\label{rlc:eq:stieltjesprimitive}
 \boxed{\displaystyle
 \int_a^b\gamma_n(c)\dd c
 =\frac{(-1)^n}{n+1}
       [\zeta^{[n+1]}(0,a)-\zeta^{[n+1]}(0,b)].}
\end{equation}
Equivalently an indefinite primitive is
$(-1)^{n+1}\zeta^{[n+1]}(0,a)/(n+1)$ plus a constant.
These are the classical Hurwitz order-derivative primitives, now attached
to the ordinary product identities of Theorem~\ref{rlc:thm:stieltjes}.
We do not count the familiar $n=0$ log-Gamma formula as a new result.

The odd-depth removable primitive also has a finite Gamma value:
\[
 \eta^{[1]}(0,a)=\log\Gamma(a/2)-\log\Gamma((a+1)/2)
                         -\tfrac12\log2.
\]
This follows from \eqref{rlc:eq:eta} and the classical value
$\zeta^{[1]}(0,a)$. In particular, the only potentially singular
coefficient in \eqref{rlc:eq:oddprimitive} is resolved explicitly by
Gamma differences rather than an unspecified integration constant.

\section{Audit findings, verification, and integration status}
\label{rlc:sec:audit}

\subsection{Mathematical status and source overlap}
No new false equation was identified in the sampled canonical claims or
the incoming results used as context. Accordingly this delivery does
not manufacture an erratum. It contributes the proved reciprocal
identities above and records the exact normalization errors that would
arise from certain invalid simplifications.

The incoming bilinear article's same-direction integer-order master
\cite{BMH} and the present all-complex-order reciprocal master are
complementary. The former's arbitrary-order same-direction and
minimal-depth questions are not solved merely by replacing its kernel
with ours. The unequal-twist article \cite{UTJ} already establishes
substantial periodic mixed-jet results; our elementary rigidity and
partial-fraction arguments should be understood as their ordinary
real-line counterparts, not as a new general monodromy theorem.
The nonlinear-transport article \cite{NST} concerns coordinate-dependent
finite parts, whereas no finite-part transport is used in the present
convergent convolution family.

\begin{center}
\begin{tabular}{@{}p{0.54\textwidth}p{0.39\textwidth}@{}}
\toprule
Claim & Status in this delivery\\
\midrule
Reciprocal all-depth, all-complex-order closure & Proved, with weighted
holomorphy and finite formulas.\\[3pt]
Every mixed logarithmic moment & Proved by a terminating polynomial
recurrence; two-factor closed form supplied.\\[3pt]
Stieltjes products, harmonic lattice collapses, and anchored primitives
& Proved exact identities.\\[3pt]
Pure order-addition rigidity for unequal shifts & Proved within the
specified functional family.\\[3pt]
Canonical Gaussian $S_4$ & Already proved in the source; unchanged.\\[3pt]
Canonical Gaussian $S_6$ and revised $S_8$ & Not proved or disproved here.\\[3pt]
Historical priority for every displayed identity & Not claimed.\\[3pt]
Proof-assistant or interval certification & Not supplied.\\
\bottomrule
\end{tabular}
\end{center}

\subsection{Concrete safeguards and corrected intermediate work}
The following distinctions are necessary for a correct integration.

\paragraph{Retain entire resonance products.}
Use $\E_a(v)$ and the even-depth form
\eqref{rlc:eq:evencentral}, rather than evaluating a vanishing
Pochhammer symbol separately from a zeta pole. Formula
\eqref{rlc:eq:resonance} loses $1/3$ under that invalid operation.

\paragraph{Retain the first-moment factor.}
The correct first moment is $(s-t)\E_a(s+t+1)/2$.
An intermediate calculation during this work omitted the $1/2$.
The independent transform argument and the tilted-series coefficient
both detect the error. It has been corrected in the final formulas and
is retained as a deliberately failing negative control in the code.
This was an error in intermediate work, not an error attributed to the
repository.

\paragraph{Do not differentiate an integer-only theorem in its integer index.}
The unequal-shift partial fraction formula holds for specified positive
integer orders. Its finite number of terms does not itself extend
holomorphically when those integers are replaced by independent complex
orders. Use \eqref{rlc:eq:unequaltransform} or the appropriate Dirichlet
mixture for that operation. In contrast, the common-shift moment formula
\eqref{rlc:eq:momentclosed} is genuinely entire in its independent orders.

\paragraph{Keep branch-cancelled sums combined.}
At even depth and positive $\mu$, the individual Lerch terms in
\eqref{rlc:eq:lerchclosure} lie on a cut. Their jump cancels because
$p_r(0)=0$. A separated branch assignment that ignores that calculation
is not justified. At $0<\Rea a\le1/2$, use an admissible nonsymmetric
vertical line rather than the Gamma-weighted line $p=1/2+it$.

\paragraph{Keep remainders and derivative types distinct.}
The product in \eqref{rlc:eq:remainderidentity} is integrable even when
its separately expanded pieces are not. Also, $\partial_s\Li_s$,
$\partial_y\Li_s$, and $\Ima\Li_s$ are different operations.
The analytic Gaussian odd part \eqref{rlc:eq:Adef} is required at
complex order.

\subsection{Executed computational checks}
%% Generated from completed verification.json.
\newcommand{\ExactCount}{336}
\newcommand{\NumericCount}{91}
\newcommand{\NegativeCount}{3}
\newcommand{\WorkingDigits}{55}

The delivered script \src{code/verify.py} executed \ExactCount\ exact
symbolic assertions, \NumericCount\ floating-point diagnostics at
\WorkingDigits\ working decimal digits, and \NegativeCount\ negative
controls. The exact checks use rational polynomial algebra, with
$P=\pi^2$ treated as a formal variable. The numerical diagnostics use
independent contour and real-integral evaluations, not fitted integer
relations.

\begin{center}
\begin{tabular}{@{}lr@{}}
\toprule
Exact assertion family & Count\\
\midrule
Central products, parity, and cosecant differential equation &36\\
Two-factor pole-cancellation and reciprocity polynomials &32\\
Weighted-transform differential polynomials &9\\
Stirling versus Bell--harmonic derivatives &136\\
Unequal-shift partial fraction examples &4\\
Negative-order residue term &1\\
First-moment conservation &11\\
Independent all-depth versus two-factor moment reductions &25\\
Consecutive-shift complete-harmonic coefficients &80\\
Explicit triple and quadruple lattice values &2\\
\midrule
Total exact assertions &336\\
\bottomrule
\end{tabular}
\end{center}

The numerical families include complex and negative orders at depths
one through eight; shifts below $1/2$ on nonsymmetric contours;
direct convolutions of elementary logistic kernels; direct reciprocal
polylogarithm moments, including a noninteger-order example; nonzero
fugacity on both sides of the apparent positive cut; Stieltjes and
log-Gamma integrals; the two unequal-shift lattice examples; and the
arctangent and first harmonic-remainder integrals.

\textbf{Evidentiary limit.} These are floating-point diagnostics, not
rigorous enclosures. Working precision is not an asserted error bound.
The harmonic-remainder quadrature uses a finite logarithmic cutoff,
which is recorded in the script. The analytic proofs, not the small
residuals, establish the infinite and continuum assertions. The complete
machine-readable results, thresholds, package versions, and generated
coefficient tables are included in \src{data/}.

\subsection{Integration recommendation}
The package is additive. A natural home beneath the Mellin report is
\begin{center}\small
\src{continuations/reciprocal-lerch-convolutions/}
\end{center}
or a dedicated report path chosen by the maintainer.
All theorem and equation labels have the prefix \src{rlc:}.
The integration note contains a short proposed manuscript insertion,
a theorem index, the source pin, and explicit conjecture-status wording.
Neither the repository nor the user's persistent Library was changed.

\section{Further research questions and topics}
\label{rlc:sec:research}

The following questions seek exact identities or structural
classifications. They are not claims inferred from small numerical
samples. The first four are especially direct continuations of the
proved formulas.

\begin{enumerate}[label=\textbf{Q\arabic*.},leftmargin=2.8em,itemsep=8pt]
\item \textbf{A full multivariate tilt formula.}
Theorem~\ref{rlc:thm:allmoments} evaluates every Taylor coefficient of
$\int e^{\sum_j\lambda_jx_j}\prod_jf_{a,s_j}(x_j)$ under
$\sum x_j=\mu$. Find a compact analytic generating function in all
$\lambda_j$, with its collision limits explicit. The two-factor
series \eqref{rlc:eq:tilted} and the moment-conservation identity supply
nonnegotiable normalization checks.

\item \textbf{Exact unequal-shift mixed order jets.}
Use \eqref{rlc:eq:unequaltransform} to construct a normally convergent
expansion into one central shift, including every independent order
derivative. The periodic analogue already exists in the incoming
unequal-twist report; the real-line version should state its own
admissible contour, convergence domain, and continuation across shift
collisions. Identify which special combinations still have finite
closure without confusing that question with integer partial fractions.

\item \textbf{Odd integral-shift cancellation sets.}
For odd $r$, the common alternating-zeta coefficient in the unit-order
lattice calculation is
$S(\boldsymbol N)=\sum_\nu c_\nu(-1)^{N_\nu}$.
Classify the distinct integer shift sets for which this rational number
is zero. Equal parity is a proved sufficient condition. A classification
would produce additional elementary multivariate integrals in
$\mathbb Q[\pi^2]$. Conversely, $S\ne0$ alone would not prove arithmetic
independence of the surviving odd-zeta combination.

\item \textbf{A central-factorial normal form for moment identities.}
Determine the exact linear relations, as functions of $W,a,\mu$, among
the finite families $\CC_{r+2q,a}(W+d;\mu)$ produced by the moment
algorithm. The product recurrence and covariant ladder already create
relations. A canonical finite reduction would improve symbolic output
without invoking conjectural independence of special values.

\item \textbf{Commensurate dilations of logarithmic arguments.}
Replace profiles by $f_{a_j,s_j}(q_jx)$ for positive rational $q_j$.
Their transforms have cosecant factors with different periods.
Investigate whether a common-denominator decomposition and the Gamma
multiplication formula yield finite Hurwitz families at integer orders,
and determine which all-complex-order sectors retain a finite closure.
The equal-dilation result here is the required base case.

\item \textbf{Complex shift boundaries and explicit contour contacts.}
Continue $a$ beyond the right half-plane while tracking the resolvent
branch point and the poles of the cosecant factor. Derive exact residue
and branch-integral corrections when a chosen contour is crossed.
The ordinary formula for $\Rea a>0$ must remain separate from any new
finite-part convention at the boundary.

\item \textbf{Higher anchored primitives at unequal collisions.}
Construct simultaneous primitives in several independent shifts, with
all constants fixed on a chosen base point and with confluent limits
stated before partial-fraction simplification. The commutator
$[\partial_a,J_a]=-J_a^2$ gives a fully proved one-shift model; the
multishift problem should identify the corresponding mixed correction
operators rather than declare them to commute.

\item \textbf{The same-direction Mellin problem beyond integer orders.}
Return to \eqref{rlc:eq:priorproblem}. Determine a spectral representation
that specializes to the incoming finite double-Hurwitz master but is
holomorphic in both orders. The reciprocal geometry provides a useful
solvable comparison, not a proof that the same finite basis works in the
same-direction geometry.

\item \textbf{The missing Gaussian functional relation.}
Seek a convergent higher-depth or octahedral relation proving the
canonical $S_6$ vector and then the revised $S_8$ vector. The formal
single-product quotient in the manuscript already supplies a separating
functional, so enlarging only that same row span cannot work. A useful
connection to the present Gaussian integral family would have to
produce the missing mixed-color relation explicitly.

\item \textbf{Noninteger convolution depth.}
The contour expression $g(p)^\rho(p+a-1)^{-W}$ can be studied for
noninteger $\rho$ after a branch is fixed. Determine the appropriate
real kernel and exact special-function expression, and identify
precisely why the finite central-factorial product is special to
integer depth. No finite-closure impossibility should be asserted
without specifying the proposed target function class.

\item \textbf{Proof-assistant decomposition and literature reconciliation.}
Formalize the polynomial recurrences, divisibility of $Q_k$, harmonic
coefficient identities, and lattice partial fractions independently of
the analysis. Then isolate weighted Banach-space holomorphy and Fourier
inversion as analytic obligations. Before assigning priority labels,
compare the exact statement set with the three incoming packages whose
contents were not inspected and with the classical convolution
literature.
\end{enumerate}

\section*{Conclusion}
\addcontentsline{toc}{section}{Conclusion}
The reciprocal constraint exposes a total-order multiplier whose
cosecant and resolvent factors yield central-factorial kernels and
Hurwitz order addition. Their interaction gives a finite calculus at
every integer depth and every complex spectral order, including all
logarithmic moments. Stieltjes products, harmonic remainders, Gamma
weights, arithmetic shift collapses, and anchored primitives follow
with their normalization and resonance terms intact.

\begin{thebibliography}{99}
\addcontentsline{toc}{section}{References}
\bibitem{Repo}
V. Reshetnikov, \emph{ProveIt: polylogarithm manuscript}, canonical manuscript, repository snapshot
\pin. Inspected files include
\src{Analysis/Polylogarithms/docs/manuscript/README.md} and
\src{chapters/04-cyclotomic-quotients.tex}.
\url{https://github.com/VladimirReshetnikov/ProveIt}.

\bibitem{MLD}
ProveIt research continuation, the \emph{Mellin--Lerch dilation continuation}, especially
\src{sections/06-audit-research.tex}, at the same source pin.
Repository path:
\src{Analysis/Polylogarithms/docs/reports/stieltjes-correlation-calculus/continuations/mellin-lerch-dilation/}.

\bibitem{BMH}
Research continuation prepared for ProveIt,
\emph{Bilinear Mellin--Hurwitz Calculus: Double-zeta reductions, harmonic
moments, and Stieltjes reciprocity}, 10 October 2026.
Incoming package \src{ProveIt_Bilinear_Mellin_Hurwitz_2026-10-10.zip};
relevant source sections 07--09 inspected from the Library copy.

\bibitem{UTJ}
Research continuation prepared for ProveIt,
\emph{Unequal-Twist Stieltjes Jets: Monodromy Obstructions, Convergent
Harmonic Expansions, Gamma Convolutions, and Multishift Zeta Finite Parts},
10 October 2026. Incoming package
\src{ProveIt_Unequal_Twist_Jets_2026-10-10.zip}; relevant source
statements and scope inspected from the Library copy.

\bibitem{NST}
Research continuation prepared for ProveIt,
\emph{Nonlinear Stieltjes Transport: Finite-part cocycles, harmonic contact
terms, M\"obius averages, and exact polylogarithmic identities},
10 October 2026. Incoming package
\src{ProveIt_Nonlinear_Stieltjes_Transport_2026-10-10.zip};
abstract and source scope inspected.

\bibitem{DLMF}
NIST, \emph{Digital Library of Mathematical Functions},
\S\S5.5, 5.12, 25.2, 25.11, 25.12, and 25.14, accessed
10 October 2026. In particular: Euler reflection and beta integrals;
Hurwitz shift and argument-derivative identities;
$\zeta'(0,a)$; the Lerch integral; and the Mellin--Barnes representation
25.14.8.
\url{https://dlmf.nist.gov/}.

\bibitem{Selvaggi}
J. P. Selvaggi and J. A. Selvaggi,
``The Application of Real Convolution for Analytically Evaluating
Fermi--Dirac-Type and Bose--Einstein-Type Integrals,''
\emph{Journal of Complex Analysis} \textbf{2018} (2018),
Article ID 5941485, 8 pages.
\url{https://doi.org/10.1155/2018/5941485}.
The convolution and Laplace-inversion methodology is an antecedent;
no claim of exhaustive overlap analysis is made.
\end{thebibliography}

\end{document}
